The trace checkpoints' shorter responses are accompanied by a numerical accuracy advantage at the 2K limit, where the unchanged student often fails to terminate. At 4K and 8K its numerical point estimate is higher. This pattern, together with the original-output diagnostic, supports completion behavior as an important part of the short-limit result. It does not identify how a truncated response would have answered if continued.

Long supervision also coexists with shorter student outputs. The serialized trace-arm targets average 5270 supervised tokens per example; at 8K the students generate 1520 tokens on average, versus 2012 unchanged. The training and test tasks differ, and teacher length was not varied. These observations do not establish that the student learned to compress the teacher's reasoning.

The no-thinking controls change the interpretation of the failed answer-only arm. Switching off the native thinking channel is insufficient to reproduce its collapse. They also provide a practical baseline for short-output use: the sampled native policy's observed accuracy is close to the 2K trace mean and its responses are much shorter. The policy comparison remains descriptive; token counts alone do not measure latency or proportional GPU-cost savings.
